\subsubsection{Three generators of degrees \texorpdfstring{$(2,2,g)$}{(2,2,g)}}\label{sec:two-block-extension}

Quadratic Jordan closure \cite[Proposition~37(b)]{MendesRadeschi}
and Euler-operator bigrading \cite[Proposition~19(a)]{MendesRadeschiQuadratic}
give a two-block, separately harmonic normal form for Laplacian
algebras generated by two quadratics and one higher-degree polynomial.
This includes the OT--FKM splitting quartics; higher-degree examples
appear at the section's end.

\begin{proposition}\label{prop:two-block-normal}
Let $\cA=\R[R,D,F]$ be a graded Laplacian algebra with independent
quadratic generators $R,D$ and a homogeneous generator $F$ of degree
$e>2$, with $F\notin\R[R,D]$. After an orthogonal
change of coordinates and a linear change of $R,D$, its quadratic
part is generated by $u=|x|^2$, $v=|y|^2$, where
$x\in\R^p$, $y\in\R^q$. There is a radial polynomial $H(u,v)$
such that $P=F-H$ is bihomogeneous of some bidegree $(a,b)$,
$a+b=e$, and $\Delta_xP=\Delta_yP=0$.

Suppose that $p,q\geq2$ and both bidegrees are positive. Then closure
is equivalent to identities
\begin{align}
 |\nabla_xP|^2&=a^2\bigl(Ku^{a-1}v^b+
               Lu^{a/2-1}v^{b/2}P\bigr),\label{eq:two-block-x}\\
 |\nabla_yP|^2&=b^2\bigl(Ku^av^{b-1}+
               Lu^{a/2}v^{b/2-1}P\bigr),\label{eq:two-block-y}
\end{align}
for constants $K>0$ and $L\in\R$. The terms involving $L$ are
omitted unless $a,b$ are both even. Necessarily
\begin{equation}\label{eq:two-block-dimensions}
 \frac{p-2}{a}=\frac{q-2}{b}.
\end{equation}
In particular, $P$ restricts to an isoparametric function on
$\Sph^{p-1}\times\Sph^{q-1}$.
\end{proposition}

\begin{proof}
Write $D(z)=\langle Bz,z\rangle$, with $B$ symmetric. The product
rule gives $\Gamma(D,D)\in\cA_2=\operatorname{span}\{R,D\}$,
so $B^2$ is a linear combination of $B$ and the identity. Since
$D$ is independent of $R$, $B$ has exactly two eigenvalues.
Its orthogonal eigenspaces give $u,v$.

Put $\mathcal B=\R[u,v]$. The operator
$x\cdot\nabla_x=\Gamma(u,\cdot)/2$ preserves $\cA$, and
$\cA_e=\mathcal B_e+\R F$. Decomposing $F$ by bidegree therefore
shows that all its nonradial terms have the same $x$-degree $a$.
Subtract the other radial terms to obtain a bihomogeneous generator
$P_0$ of bidegree $(a,b)$. This algebra is now bigraded. The distinct
bidegree components of $\Delta P_0$ belong to $\mathcal B$, since
they have degree less than $e$.

They vanish unless $a,b$ are even. In the even mixed case write
$\Delta_xP_0=\lambda_xu^{a/2-1}v^{b/2}$ and
$\Delta_yP_0=\lambda_yu^{a/2}v^{b/2-1}$. Commuting these partial
Laplacians gives
$\lambda_x b(b+q-2)=\lambda_y a(a+p-2)$.
Subtracting the appropriate multiple of $u^{a/2}v^{b/2}$ kills
both. If one bidegree is zero, the subtraction only uses the active
block. This proves the separately harmonic normal form.

The two bidegree components of $\Gamma(P,P)$ also lie in the
algebra. Their degree is $2e-2$, so their expressions contain at
most the first power of $P$. Matching bidegrees gives
\eqref{eq:two-block-x}--\eqref{eq:two-block-y} initially with
possibly different constant and linear coefficients. On the unit
product spheres, let $f=P|_{\Sph^{p-1}\times\Sph^{q-1}}$.
The separate tangential gradient squares are quadratics in $f$,
with leading coefficients $-a^2$ and $-b^2$. Its mean is zero,
by separate harmonicity. Its maximum $M>0$ and minimum $m<0$
are roots of both quadratics, since both partial gradients vanish
at these points. Thus the two quadratics are respectively
$a^2(K+Lf-f^2)$ and $b^2(K+Lf-f^2)$, with
$K=-mM>0$ and $L=m+M$. This proves the common coefficients.
Conversely, the two identities and the separate Euler and Laplace
identities give every generator entry of the closure table, so
\cref{prop:MR-generator-test} proves sufficiency.

If $\mu_2$ is the average of $f^2$, integration by parts on each
sphere gives
\[
 a^2K=a(2a+p-2)\mu_2,\qquad
 b^2K=b(2b+q-2)\mu_2.
\]
This implies \eqref{eq:two-block-dimensions}. Finally,
\[
 |\nabla f|^2=(a^2+b^2)(K+Lf-f^2),\qquad
 \Delta f=-[a(a+p-2)+b(b+q-2)]f,
\]
which are the isoparametric equations on the product.
\end{proof}

The block dimensions $p,q$ need
not agree; if $p=q>2$, \eqref{eq:two-block-dimensions} forces $a=b$.
The change $F\mapsto F-H(u,v)$ preserves the algebra but generally
changes its zero set, so the minimal level must still be determined
within each explicit family. If one bidegree is zero, the active block
reduces to \cite[Proposition~37(c)]{MendesRadeschi}, with an independent
Euclidean factor adjoined.

